% Part DetPoly: one chapter per topic of `DetPoly/`.
% Throughout, R is a commutative ring and n a finite index type.

\chapter{Determinantal polynomials}

\begin{definition}[Submatrices and principal minors]
\label{def:principal-minor}
\alsoIn{part:warmup}
\lean{Matrix.setSubmatrix, Matrix.principalMinor, Matrix.submatrix_sumCompl_eq_fromBlocks}
\leanok
For $S,T\subseteq n$, $A[S,T]$ is the submatrix on rows $S$ and columns $T$.
For $S\subseteq n$ finite, $A_S=\det A[S,S]$.
\end{definition}

\begin{definition}[Determinantal polynomial]
\label{def:detpoly}
\alsoIn{part:warmup}
\lean{Matrix.mvCharmatrix, Matrix.detpoly, Matrix.mvCharmatrix_submatrix}
\leanok
$M_A=\operatorname{diag}(x)+A\in R[x]^{n\times n}$ and $f_A=\det M_A$.
\end{definition}

\begin{definition}[Rayleigh difference]
\label{def:rayleigh}
\alsoIn{part:warmup}
\lean{rayleighDiff}
\leanok
$\Delta_{ij}(f)=\partial_i f\,\partial_j f-f\,\partial_i\partial_j f$.
\end{definition}

\begin{theorem}[Principal minor expansion]
\label{thm:det-diagonal-add}
\lean{Matrix.det_diagonal_add}
\leanok
\uses{def:principal-minor}
For $x\in R^n$, $\det(\operatorname{diag}(x)+A)=\sum_{S\subseteq n}\prod_{i\in S}x_i\cdot A_{S^c}$.
\end{theorem}

\begin{proof}
\leanok
Expand by multilinearity in the rows, choosing the row $x_ie_i$ for $i\in S$ and the row
$A_i$ otherwise.
\end{proof}

\begin{corollary}[Coefficients of $f_A$]
\label{cor:coeff-detpoly}
\alsoIn{part:warmup}
\lean{Matrix.detpoly_eq_sum, Finset.sum_single_one_injective, Matrix.coeff_detpoly,
  Matrix.detpoly_mem_restrictDegree, detpoly_eq_iff_principalMinor_eq}
\leanok
\uses{def:detpoly}
The coefficient of $\prod_{i\in S}x_i$ in $f_A$ is $A_{S^c}$. Hence $f_A$ is multiaffine, and
$f_A=f_B$ iff $A_S=B_S$ for all $S$.
\end{corollary}

\begin{proof}
\leanok
\uses{thm:det-diagonal-add}
Apply Theorem~\ref{thm:det-diagonal-add} over $R[x]$; distinct $S$ give distinct monomials.
\end{proof}

\begin{theorem}[Desnanot--Jacobi]
\label{thm:desnanot-jacobi}
\alsoIn{part:warmup}
\lean{Matrix.det_toSquareBlockProp_pair, Matrix.adjugate_mul_adjugate_sub_adjugate_mul_adjugate}
\leanok
For $M\in R^{n\times n}$ and $i\ne j$,
$\operatorname{adj}(M)_{ii}\operatorname{adj}(M)_{jj}-\operatorname{adj}(M)_{ij}\operatorname{adj}(M)_{ji}
=\det M\cdot\det M[\{i,j\}^c,\{i,j\}^c]$.
\end{theorem}

\begin{proof}
\leanok
Let $C$ be the identity with columns $i,j$ replaced by those of $\operatorname{adj}M$. Both $C$
and $MC$ are block triangular along $\{i,j\}^c$, which gives
$\det M\cdot\mathrm{LHS}=\det M\cdot\mathrm{RHS}$. Cancel $\det M$ for the generic matrix over
$\mathbb Z[x_{kl}]$ and specialize.
\end{proof}

\begin{lemma}[Derivatives of $f_A$]
\label{lem:pderiv-detpoly}
\alsoIn{part:warmup}
\lean{pderiv_detPoly_eq, pderiv_pderiv_detPoly_eq}
\leanok
\uses{def:detpoly}
$\partial_i f_A=\det M_A[\{i\}^c,\{i\}^c]$, and for $i\ne j$,
$\partial_i\partial_j f_A=\det M_A[\{i,j\}^c,\{i,j\}^c]$.
\end{lemma}

\begin{proof}
\uses{cor:coeff-detpoly}
$x_i$ occurs only in the $(i,i)$ entry; expand along row $i$, or compare the expansions of
Corollary~\ref{cor:coeff-detpoly} for $A$ and $A[\{i\}^c,\{i\}^c]$.
\end{proof}

\begin{theorem}[Al Ahmadieh, Thm 2.3]
\label{thm:rayleigh-detpoly}
\alsoIn{part:warmup}
\lean{theorem2_3}
\leanok
\uses{def:rayleigh, def:detpoly}
With $G=\operatorname{adj}M_A$, $\Delta_{ij}(f_A)=G_{ij}G_{ji}$.
\end{theorem}

\begin{proof}
\uses{lem:pderiv-detpoly, thm:desnanot-jacobi}
For $i\ne j$, $\partial_if_A=G_{ii}$ and $\partial_i\partial_jf_A=\det M_A[\{i,j\}^c,\{i,j\}^c]$;
apply Theorem~\ref{thm:desnanot-jacobi} to $M_A$.
\end{proof}

\chapter{Diagonal equivalence and irreducibility}

\begin{definition}[Diagonal similarity and equivalence]
\label{def:diag}
\alsoIn{part:warmup}
\lean{Matrix.DiagSimilar, Matrix.DiagEquivalent, Matrix.DiagSimilar.equivalence,
  Matrix.DiagEquivalent.equivalence, Matrix.diagSimilarSetoid, Matrix.diagEquivalentSetoid,
  Matrix.DiagSimilar.transpose, Matrix.DiagSimilar.diagEquivalent,
  Matrix.diagEquivalent_iff_diagSimilar_of_isSymm}
\leanok
$M\sim_{DS}N$ if $M=DND^{-1}$ for an invertible diagonal $D$.
$M\sim_{DE}N$ if $M\sim_{DS}N$ or $M\sim_{DS}N^T$. Both are equivalence relations.
\end{definition}

\begin{lemma}[Diagonal similarity preserves zero patterns]
\label{lem:diag-zero}
\lean{Matrix.DiagSimilar.apply_eq_zero_iff}
\leanok
\uses{def:diag}
If $M\sim_{DS}N$ then $M_{ij}=0\iff N_{ij}=0$.
\end{lemma}

\begin{proof}
\leanok
$M_{ij}=d_iN_{ij}d_j^{-1}$.
\end{proof}

\begin{lemma}[Diagonal equivalence preserves principal minors]
\label{lem:diag-principal-minor}
\lean{Matrix.principalMinor_transpose, Matrix.DiagSimilar.principalMinor_eq,
  Matrix.DiagEquivalent.principalMinor_eq}
\leanok
\uses{def:diag, def:principal-minor}
If $M\sim_{DE}N$ then $M_S=N_S$ for all $S$.
\end{lemma}

\begin{proof}
\leanok
$(DND^{-1})[S,S]=D_SN[S,S]D_S^{-1}$, and $\det N^T=\det N$.
\end{proof}

\begin{definition}[Irreducible matrix]
\label{def:irred}
\alsoIn{part:warmup}
\lean{Matrix.toQuiver', Matrix.IsIrred, Matrix.IsReducible}
\leanok
$A$ is irreducible if the digraph with an edge $i\to j$ iff $A_{ij}\ne0$ is strongly
connected, and reducible otherwise.
\end{definition}

\begin{lemma}[Reducibility via a set]
\label{lem:reducible-iff}
\alsoIn{part:warmup}
\lean{Matrix.isReducible_iff_exists_set, Matrix.isReducible_iff_exists_finset,
  Matrix.isReducible_iff_blockTriangular, Matrix.blockTriangular_of_subset}
\leanok
\uses{def:irred}
$A$ is reducible iff there is a nonempty proper $S\subseteq n$ with $A_{ij}=0$ for all
$i\notin S$, $j\in S$, iff $A$ is block triangular for a surjective labelling $n\to\{0,1\}$.
\end{lemma}

\begin{proof}
\leanok
Take $S$ to be the set of vertices that can reach a vertex $j$ not reachable from some $i$.
\end{proof}

\begin{lemma}[Determinants of block triangular matrices]
\label{lem:blocktri-det}
\lean{Matrix.BlockTriangular.det_eq_det}
\leanok
Two matrices that are block triangular for the same labelling and agree on the diagonal
blocks have the same determinant.
\end{lemma}

\begin{proof}
\leanok
Both determinants are the product of the determinants of the diagonal blocks.
\end{proof}

\begin{theorem}[Al Ahmadieh, Thm 4.1 ($\Rightarrow$), reducible case]
\label{thm:reducible-fiber}
\alsoIn{part:warmup}
\lean{Matrix.IsReducible.exists_principalMinor_eq_not_diagEquivalent}
\leanok
\uses{def:irred, def:diag, def:principal-minor}
If $R$ is nontrivial and $A$ is reducible, some $B$ has the same principal minors as $A$ but
$B\not\sim_{DE}A$.
\end{theorem}

\begin{proof}
\leanok
\uses{lem:reducible-iff, lem:blocktri-det, lem:diag-zero}
Write $A=\begin{pmatrix}E&F\\0&G\end{pmatrix}$ and flip whether one entry of $F$ is zero.
Principal minors only see $E$ and $G$, while diagonal equivalence preserves the zero pattern
of $A$ or of $A^T$.
\end{proof}

\chapter{Rank-one updates}

\begin{lemma}[Matrix determinant lemma]
\label{lem:det-add-vecMulVec}
\lean{Matrix.det_add_vecMulVec}
\leanok
$\det(A+uv^T)=\det A+v^T\operatorname{adj}(A)u$.
\end{lemma}

\begin{proof}
\leanok
Add $v_ju$ to the columns $j$ one at a time; each step adds $v_j\det(A\text{ with column }j
\text{ replaced by }u)$.
\end{proof}

\begin{lemma}[Adjugate of a block matrix with a rank-one block]
\label{lem:adj-fromBlocks}
\alsoIn{part:warmup}
\lean{Matrix.adjugate_fromBlocks_vecMulVec_apply₁₂}
\leanok
For $M=\begin{pmatrix}A&uv^T\\C&D\end{pmatrix}$, $i$ in the first block and $j$ in the second,
$\operatorname{adj}(M)_{ij}=-\det(A\text{ with column }i\text{ replaced by }u)\cdot
\det(D\text{ with row }j\text{ replaced by }v)$.
\end{lemma}

\begin{proof}
\leanok
$\operatorname{adj}(M)_{ij}=\det N$ where row $j$ of $N$ is $e_i$. Column operations using that
row turn column $i$ into $(u,e_j)$ and then clear the upper right block.
\end{proof}

\begin{lemma}[Determinant of a block matrix with rank-one off-diagonal blocks]
\label{lem:det-fromBlocks}
\lean{Matrix.det_fromBlocks_vecMulVec}
\leanok
$\det\begin{pmatrix}A&pq^T\\uv^T&D\end{pmatrix}
=\det A\det D-(v^T\operatorname{adj}(A)p)(q^T\operatorname{adj}(D)u)$.
\end{lemma}

\begin{proof}
\leanok
\uses{lem:det-add-vecMulVec, lem:adj-fromBlocks}
The matrix is $\begin{pmatrix}A&pq^T\\0&D\end{pmatrix}+(0,u)(v,0)^T$. Apply
Lemma~\ref{lem:det-add-vecMulVec} and evaluate the adjugate entries in the upper right block
by Lemma~\ref{lem:adj-fromBlocks}.
\end{proof}

\begin{lemma}[Rank at most one over a field]
\label{lem:rank-le-one}
\lean{Matrix.rank_le_one_iff_eq_vecMulVec, Matrix.one_le_rank_of_apply_ne_zero}
\leanok
Over a field, $\operatorname{rk}M\le1$ iff $M=uv^T$ for some $u,v$; and
$\operatorname{rk}M\ge1$ if $M$ has a nonzero entry.
\end{lemma}

\begin{proof}
\leanok
The column space is spanned by a single vector $u$; write column $j$ as $v_ju$.
\end{proof}

\chapter{Cuts and cut-transposes}

\begin{definition}[Rank-one factorization]
\label{def:r1fact}
\lean{Matrix.R1Factorization, Matrix.R1Factorization.block₁₁, Matrix.R1Factorization.block₁₂,
  Matrix.R1Factorization.block₂₁, Matrix.R1Factorization.block₂₂,
  Matrix.R1Factorization.block₁₂_eq, Matrix.R1Factorization.block₂₁_eq,
  Matrix.R1Factorization.submatrix_sumCompl, Matrix.R1Factorization.compl,
  Matrix.R1Factorization.submatrix}
\leanok
\uses{def:principal-minor}
A set $X\subseteq n$ with vectors $p,q,u,v$ such that $A[X,X^c]=pq^T$ and $A[X^c,X]=uv^T$,
i.e.\ $A=\begin{pmatrix}E&pq^T\\uv^T&N\end{pmatrix}$ along $X,X^c$. It restricts to principal
submatrices and can be read along $X^c$.
\end{definition}

\begin{definition}[Cut-transpose, CGGR Def 2.10]
\label{def:cut-transpose}
\lean{Matrix.cutTranspose, Matrix.cutTranspose_apply, Matrix.R1Factorization.swap,
  Matrix.cutTranspose_swap, Matrix.submatrix_cutTranspose,
  Matrix.submatrix_sumCompl_cutTranspose}
\leanok
\uses{def:r1fact}
$\operatorname{ct}(A,X)=\begin{pmatrix}E&pu^T\\qv^T&N^T\end{pmatrix}$. It commutes with taking
principal submatrices, and $\operatorname{ct}(\operatorname{ct}(A,X),X)=A$ for the swapped
factorization.
\end{definition}

\begin{definition}[Cut]
\label{def:cut}
\alsoIn{part:warmup}
\lean{Matrix.R1CutFactorization, Matrix.HasCut, Matrix.hasCut_iff_exists_r1Factorization,
  Matrix.HasCut.compl, Matrix.hasCut_compl_iff, Matrix.HasCut.cutTranspose,
  Matrix.IsCutTranspose, Matrix.HasCut.exists_isCutTranspose, Matrix.IsCutTranspose.symm,
  Matrix.IsCutTranspose.hasCut_left, Matrix.IsCutTranspose.hasCut_right}
\leanok
\uses{def:cut-transpose}
$X$ is a cut of $A$ if $|X|\ge2$, $|X^c|\ge2$ and $A$ has a rank-one factorization along $X$.
$B$ is a cut-transpose of $A$ at $X$ if $B=\operatorname{ct}(A,X)$ for such a factorization.
\end{definition}

\begin{definition}[Cut-transpose paths]
\label{def:ct-path}
\lean{Matrix.CutTransposePath, Matrix.CutTransposePath.length, Matrix.CutTransposePath.cuts,
  Matrix.CutTransposePath.append, Matrix.CutTransposePath.reverse,
  Matrix.CutTransposeConnected, Matrix.cutTransposeConnected_equivalence,
  Matrix.IsCutTranspose.exists_path}
\leanok
\uses{def:cut}
A sequence $A=A_0,\dots,A_k=C$ with $A_{i+1}=\operatorname{ct}(A_i,X_i)$ for cuts $X_i$ of
$A_i$. Being connected by such a path is an equivalence relation.
\end{definition}

\begin{lemma}[Transposes and complements]
\label{lem:ct-transpose}
\lean{Matrix.R1Factorization.transpose, Matrix.transpose_cutTranspose,
  Matrix.cutTranspose_compl, Matrix.HasCut.transpose, Matrix.hasCut_transpose,
  Matrix.IsCutTranspose.transpose, Matrix.IsCutTranspose.compl}
\leanok
\uses{def:cut}
$\operatorname{ct}(A,X)^T=\operatorname{ct}(A^T,X)$ and
$\operatorname{ct}(A,X^c)=\operatorname{ct}(A,X)^T$. $X$ is a cut of $A^T$ iff it is a cut of
$A$.
\end{lemma}

\begin{proof}
\leanok
Entrywise.
\end{proof}

\begin{lemma}[Cut-transpose preserves the determinant]
\label{lem:det-ct}
\lean{Matrix.det_cutTranspose}
\leanok
\uses{def:cut-transpose}
$\det\operatorname{ct}(A,X)=\det A$, for any rank-one factorization (no size condition).
\end{lemma}

\begin{proof}
\leanok
\uses{lem:det-fromBlocks}
By Lemma~\ref{lem:det-fromBlocks},
$\det A=\det E\det N-(v^T\operatorname{adj}(E)p)(q^T\operatorname{adj}(N)u)$, which is
invariant under $q\leftrightarrow u$, $N\mapsto N^T$.
\end{proof}

\begin{theorem}[CGGR Lemma 2.12]
\label{thm:pm-ct}
\lean{Matrix.principalMinor_cutTranspose, Matrix.IsCutTranspose.principalMinor_eq,
  Matrix.CutTransposePath.principalMinor_eq, Matrix.CutTransposeConnected.principalMinor_eq}
\leanok
\uses{def:cut, def:ct-path, def:principal-minor}
Cut-transposes, and hence cut-transpose paths, preserve principal minors.
\end{theorem}

\begin{proof}
\leanok
\uses{lem:det-ct}
A principal submatrix of $\operatorname{ct}(A,X)$ is the cut-transpose of the corresponding
principal submatrix of $A$; apply Lemma~\ref{lem:det-ct}.
\end{proof}

\begin{lemma}[Cuts over a field]
\label{lem:cut-rank}
\lean{Matrix.R1Factorization.rank_le_one, Matrix.exists_r1Factorization_of_rank_le_one,
  Matrix.hasCut_iff_rank_le_one, Matrix.HasCut.rank_eq_one,
  Matrix.R1CutFactorization.rank_eq_one}
\leanok
\uses{def:cut, def:irred}
Over a field, $X$ is a cut iff $|X|,|X^c|\ge2$ and
$\operatorname{rk}A[X,X^c],\operatorname{rk}A[X^c,X]\le1$. If $A$ is irreducible, both ranks
equal $1$.
\end{lemma}

\begin{proof}
\leanok
\uses{lem:rank-le-one, lem:reducible-iff}
Lemma~\ref{lem:rank-le-one}. If a cross block were zero, $X$ or $X^c$ would witness
reducibility (Lemma~\ref{lem:reducible-iff}).
\end{proof}

\chapter{Lemma 4.3}

Throughout, $\mathbb F$ is a field and $G=\operatorname{adj}M_A$ over $\mathbb F[x]$.

\begin{lemma}[Al Ahmadieh, Lemma 4.3, minors]
\label{lem:43-minors}
\lean{Matrix.hasCut_minors_multiples}
\leanok
\uses{def:cut, def:detpoly, def:irred}
If $A$ is irreducible with a cut $X$, there are $a_i\in\mathbb F[x_k:k\in X\setminus\{i\}]$
and $c_j\in\mathbb F$ such that for $i\in X$, $j\notin X$, the minor of $M_A$ on rows $X$ and
columns $X$ with $i$ replaced by $j$ equals $c_ja_i$.
\end{lemma}

\begin{proof}
Write $A[X,X^c]=pq^T$. Replacing column $i$ of $M_A[X,X]$ by column $j$ of $A[X,X^c]$, which
is $q_jp$, and pulling out $q_j$ leaves a determinant free of $x_i$ and of $j$.
\end{proof}

\begin{lemma}[Al Ahmadieh, Lemma 4.3, cut form]
\label{lem:43}
\lean{Matrix.hasCut_adjPoly_of_hasCut}
\leanok
\uses{def:cut, def:detpoly, def:irred}
If $A$ is irreducible with a cut $X$, then $X$ is a cut of $G$ over $\mathbb F[x]$.
\end{lemma}

\begin{proof}
\uses{lem:adj-fromBlocks}
Reindex $M_A$ along $X\oplus X^c$ (Definition~\ref{def:r1fact}). By
Lemma~\ref{lem:adj-fromBlocks}, $G_{ij}=-\det(E_x\text{ with column }i\text{ replaced by }p)
\cdot\det(N_x\text{ with row }j\text{ replaced by }q)$ for $i\in X$, $j\notin X$, an outer
product. Do the same for $X^c$.
\end{proof}
